\section{Triple Integration with Cylindrical and Spherical Coordinates}\label{sec:cylindrical_spherical}

Just as polar coordinates gave us a new way of describing curves in the plane, in this section we will see how \emph{cylindrical} and \emph{spherical} coordinates give us new ways of describing surfaces and regions in space.

\mtable{Illustrating the principles behind cylindrical coordinates.}{fig:cylindricalintro}{%
\myincludeasythree{
3Droll=0,
3Dortho=0.004267800133675337,
3Dc2c=0.6733480095863342 0.6255449652671814 0.39407607913017273,
3Dcoo=-12.612065315246582 -11.544181823730469 33.65996551513672,
3Droo=150.0000065104165,
3Dlights=Headlamp}{\marginparwidth}{The point (r,𝜃,z) in cylindrical coordinates is z units above the (r,𝜃) point in polar coordinates in the xy plane.}{figures/figcylindricalintro_3D}}

\subsection{Cylindrical Coordinates}

In short, cylindrical coordinates can be thought of as a combination of the polar and rectangular coordinate systems. %One can identify a location in the $x$-$y$ plane using polar coordinates (i.e., using $(r,\theta)$), then use a $z$-value to determine the distance
One can identify a point $(x_0,y_0,z_0)$, given in rectangular coordinates, with the point $(r_0,\theta_0,z_0)$, given in cylindrical coordinates, where the $z$-value in both systems is the same, and the point $(x_0,y_0)$ in the $x$-$y$ plane is identified with the polar point $P(r_0,\theta_0)$; see \autoref{fig:cylindricalintro}. So that each point in space that does not lie on the $z$-axis is defined uniquely, we will restrict $r\geq 0$ and $0\leq \theta< 2\pi$.
\index{cylindrical coordinates}\index{coordinates!cylindrical}

We use the identity $z=z$ along with the identities found in \autoref{idea:polarconvert} to convert between the  rectangular coordinate $(x,y,z)$ and the  cylindrical coordinate $(r,\theta,z)$, namely:
\begin{align*}
&\text{From rectangular to cylindrical:}\quad r=\sqrt{x^2+y^2},\quad \tan\theta = y/x \quad\text{and}\quad z=z;\\
&\text{From cylindrical to rectangular:} \quad x=r\cos\theta\quad y=r\sin\theta\quad \text{and}\quad z=z.
\end{align*}
These identities, along with conversions related to spherical coordinates, are given later in \autoref{idea:convert_coords}.

\mnote{\textbf{Note:} Our rectangular to polar conversion formulas used $r^2=x^2+y^2$, allowing for negative $r$ values. Since we now restrict $r\geq 0$, we can use $r=\sqrt{x^2+y^2}$.}

\youtubeVideo{w9wCdLuklw8}{Conversion From Rectangular Coordinates}

\begin{example}[Converting between rectangular and cylindrical coordinates]\label{ex_cylindrical4}%
Convert the rectangular point $(2,-2,1)$ to cylindrical coordinates, and convert the cylindrical point $(4,3\pi/4,5)$ to rectangular.
\solution
Following the identities given above (also in \autoref{idea:convert_coords}), we have that $r = \sqrt{2^2+(-2)^2} = 2\sqrt{2}$. Using $\tan\theta = y/x$, we find that $\theta = \tan^{-1}(-2/2) =-\pi/4$. As we restrict $\theta$ to being between $0$ and $2\pi$, we set $\theta = 7\pi/4$. Finally, $z = 1$, giving the cylindrical point $(2\sqrt2,7\pi/4,1)$.

In converting the cylindrical point $(4,3\pi/4,5)$ to rectangular, we have
$x = 4\cos\bigl(3\pi/4\bigr) = -2\sqrt{2}$, $y = 4\sin\bigl(3\pi/4\bigr) = 2\sqrt{2}$ and $z=5$, giving the rectangular point $(-2\sqrt{2},2\sqrt{2},5)$.
\end{example}

Setting each of $r$, $\theta$ and $z$ equal to a constant defines a surface in space, as illustrated in the following example.

\mtable{Graphing the canonical surfaces in cylindrical coordinates from \autoref{ex_cylindrical1}.}{fig:cylindrical1}{%
\myincludeasythree{
3Droll=0,
3Dortho=0.005137230269610882,
3Dc2c=0.8924431204795837 0.30546268820762634 0.3320206105709076,
3Dcoo=-13.9379243850708 -4.1078948974609375 30.89048194885254,
3Droo=149.9999934895832,
3Dlights=Headlamp}{\marginparwidth}{The surface z=2 is a flat plane.  The surface r=1 is a cylinder about the z axis.  The surface 𝜃=π/3 is a half plane with one edge along the z axis.}{figures/figcylindrical1_3D}}

\begin{example}[Canonical surfaces in cylindrical coordinates]\label{ex_cylindrical1}%
Describe the surfaces $r=1$, $\theta = \pi/3$ and $z=2$, given in cylindrical coordinates.
\solution
The equation $r=1$ describes all points in space that are 1 unit away from the $z$-axis. This surface is a ``tube'' or ``cylinder'' of radius 1, centered on the $z$-axis, as graphed in \autoref{fig:spacecylinder1} (which describes the cylinder $x^2+y^2=1$ in space). 

The equation $\theta=\pi/3$ describes the plane formed by extending the line $\theta=\pi/3$, as given by polar coordinates in the $x$-$y$ plane, parallel to the $z$-axis.

The equation $z=2$ describes the plane of all points in space that are 2 units above the $x$-$y$ plane. This plane is the same as the plane described by $z=2$ in rectangular coordinates.

All three surfaces are graphed in \autoref{fig:cylindrical1}. Note how their intersection uniquely defines the point $P=(1,\pi/3,2)$.
\end{example}

Cylindrical coordinates are useful when describing certain domains in space, allowing us to evaluate triple integrals over these domains more easily than if we used rectangular coordinates.

\autoref{thm:triple_integration2} shows how to evaluate $\iiint_Dh(x,y,z)\dd V$ using rectangular coordinates. In that evaluation, we use $\dd V = \dd z\dd y\dd x$ (or one of the other five orders of integration). Recall how, in this order of integration, the bounds on $y$ are ``curve to curve'' and the bounds on $x$ are ``point to point'': these bounds describe a region $R$ in the $x$-$y$ plane. We could describe $R$ using polar coordinates as done in \autoref{sec:double_int_polar}. In that section, we saw how we used $\dd A = r\dd r\dd\theta$ instead of $\dd A = \dd y\dd x$. 

Considering the above thoughts, we have $\dd V = \dd z\bigl(r\dd r\dd\theta\bigr) = r\dd z\dd r\dd\theta$. We set bounds on $z$ as ``surface to surface'' as done in the previous section, and then use ``curve to curve'' and ``point to point'' bounds on $r$ and $\theta$, respectively. Finally, using the identities given above, we change the integrand $h(x,y,z)$ to $h(r,\theta,z)$.

This process should sound plausible; the following theorem states it is truly a way of evaluating a triple integral.

\begin{theorem}[Triple Integration in Cylindrical Coordinates]\label{thm:triple_int_cylindrical}%
%
Let $w=h(r,\theta,z)$ be a continuous function on a closed, bounded region $D$ in space, bounded in cylindrical coordinates by $\alpha \leq \theta \leq \beta$, $g_1(\theta)\leq r \leq g_2(\theta)$ and $f_1(r,\theta) \leq z \leq f_2(r,\theta)$. Then \index{integration!with cylindrical coordinates} 
\[
\iiint_D h(r,\theta,z)\dd V = \int_\alpha^\beta\int_{g_1(\theta)}^{g_2(\theta)}\int_{f_1(r,\theta)}^{f_2(r,\theta)}h(r,\theta,z) r\dd z\dd r\dd\theta.
\]
\end{theorem}

\mtable{Visualizing the solid used in \autoref{ex_cylindrical2}.}{fig:cylindrical2}{%
\myincludeasythree{
3Droll=0,
3Dortho=0.004440656863152981,
3Dc2c=0.6615753173828125 0.6849026679992676 0.30533018708229065,
3Dcoo=-6.77640962600708 2.852377414703369 51.19231414794922,
3Droo=300,
3Dlights=Headlamp}{\marginparwidth}{A cylinder of radius 2 centered along the z axis where 0≤z≤3.  Once z=3, the cylinder is capped by a sphere of radius 2.}{figures/figcylindrical2_3D}}

\begin{example}[Evaluating a triple integral with cylindrical coordinates]\label{ex_cylindrical2}%
Find the mass of the solid represented by the region in space bounded by $z=0$, $z=\sqrt{4-x^2-y^2}+3$ and the cylinder $x^2+y^2=4$ (as shown in \autoref{fig:cylindrical2}), with density function $\delta(x,y,z) = x^2+y^2+z+1$, using a triple integral in cylindrical coordinates. Distances are measured in centimeters and density is measured in grams/cm$^3$.
\solution
We begin by describing this region of space with cylindrical coordinates. The plane $z=0$ is left unchanged; with the identity $r=\sqrt{x^2+y^2}$, we convert the hemisphere of radius 2 to the equation $z=\sqrt{4-r^2}$; the cylinder $x^2+y^2=4$ is converted to $r^2=4$, or, more simply, $r=2$.  We also convert the density function: $\delta(r,\theta,z) = r^2+z+1$.

To describe this solid with the bounds of a triple integral, we bound $z$ with $0\leq z\leq \sqrt{4-r^2}+3$; we bound $r$ with $0 \leq r \leq 2$; we bound $\theta$ with $0 \leq \theta \leq 2\pi$.

Using \autoref{def:mass_3d} and \autoref{thm:triple_int_cylindrical}, we have the mass of the solid is
\begin{align*}
M=\iiint_D\delta(x,y,z)\dd V &= \int_0^{2\pi}\int_0^2\int_0^{\sqrt{4-r^2}+3}\bigl(r^2+z+1\bigr)r\dd z\dd r\dd\theta \\
&= \int_0^{2\pi}\int_0^2\bigl((r^3+4r)\sqrt{4-r^2}+\frac52r^3+\frac{19}2r\bigr)\dd r\dd\theta \\
&= \frac{1318\pi}{15} \approx 276.04\text{ gm},
\end{align*}
where we leave the details of the remaining double integral to the reader.
\end{example}

\begin{example}[Finding the center of mass using cylindrical coordinates]\label{ex_cylindrical3}%
Find
%
\mtable{Visualizing the solid used in \autoref{ex_cylindrical3}.}{fig:cylindrical3}{%
\myincludeasythree{
3Droll=-0.0853586683163278,
3Dortho=0.003775296732783318,
3Dc2c=0.5381587743759155 0.7346906065940857 0.41305553913116455,
3Dcoo=-13.062450408935547 -3.628488063812256 70.90967559814453,
3Droo=300}{\marginparwidth}{Three loops joined together, with one loop tapering down to a point away from the joint.}{figures/figdoublepol4_3D}}
%
the center of mass of the solid with constant density whose base can be described by the polar curve $r=\cos(3\theta)$ and whose top is defined by the plane $z=1-x+0.1y$, where distances are measured in feet, as seen in \autoref{fig:cylindrical3}. (The volume of this solid was found in \autoref{ex_doublepol4}.)
\solution
We convert the equation of the plane to use cylindrical coordinates: $z= 1-r\cos\theta+0.1r\sin\theta$. Thus the region is space is bounded by $0 \leq z \leq 1-r\cos\theta + 0.1r\sin\theta$, $0 \leq r \leq \cos(3\theta)$, $0 \leq \theta \leq \pi$ (recall that the rose curve $r=\cos(3\theta)$ is traced out once on $[0,\pi]$.

Since density is constant, we set $\delta = 1$ and finding the mass is equivalent to finding the volume of the solid. We set up the triple integral to compute this but do not evaluate it; we leave it to the reader to confirm it evaluates to the same result found in \autoref{ex_doublepol4}.
\[
M = \iiint_D\delta \dd V = \int_0^{\pi}\int_0^{\cos(3\theta)}\int_0^{1-r\cos\theta+0.1r\sin\theta} r\dd z\dd r\dd\theta =\frac\pi4.%\approx 0.785.
\]

From \autoref{def:mass_3d} we set up the triple integrals to compute the moments about the three coordinate planes. The computation of each is left to the reader (using technology is recommended):
\begin{align*}
M_{yz} = \iiint_D x\dd V &= \int_0^{\pi}\int_0^{\cos(3\theta)}\int_0^{1-r\cos\theta+0.1r\sin\theta} (r\cos\theta) r\dd z\dd r\dd\theta\\
&= -\frac{3\pi}{64}. %-0.147.
\\
%\end{align*}
%\begin{align*}
M_{xz} = \iiint_D y\dd V &= \int_0^{\pi}\int_0^{\cos(3\theta)}\int_0^{1-r\cos\theta+0.1r\sin\theta} (r\sin\theta) r\dd z\dd r\dd\theta\\
&= \frac{3\pi}{640}. %0.015.
\\
M_{xy} = \iiint_D z\dd V &= \int_0^{\pi}\int_0^{\cos(3\theta)}\int_0^{1-r\cos\theta+0.1r\sin\theta} (z) r\dd z\dd r\dd\theta\\
 &= \frac{1903\pi}{12800}. %0.467.
\end{align*}
The center of mass %, in rectangular coordinates,  
is located at
\[(-3/16, 3/160, 1903/3200)=(-0.1875,0.01875,0.5946875),\]
which lies outside the bounds of the solid.
\end{example}

\subsection{Spherical Coordinates}
\index{spherical coordinates}\index{coordinates!spherical}

\mtable{Illustrating the principles behind spherical coordinates.}{fig:sphericalintro}{%
\myincludeasythree{
3Droll=0,
3Dortho=0.0045,
3Dc2c=0.6936245560646057 0.5807817578315735 0.4261191189289093,
3Dcoo=-15.555180549621582 -9.786354064941406 35.21320343017578,
3Droo=150,
3Dlights=Headlamp}{\marginparwidth}{The point (𝜌,𝜃,𝜑) in spherical coordinates is distance 𝜌 from the origin, and directly above the 𝜃 line from polar or cylindrical coordinates.  The line from the origin to the point makes an angle of 𝜑 with the positive z axis.}{figures/figsphericalintro_3D}}

In short, spherical coordinates can be thought of as a ``double application'' of the polar coordinate system. In spherical coordinates, a point $P$ is identified with $(\rho,\theta,\varphi)$, where $\rho$ is the distance from the origin to $P$, $\theta$ is the same angle as would be used to describe $P$ in the cylindrical coordinate system, and $\varphi$ is the angle between the positive $z$-axis and the ray from the origin to $P$; see \autoref{fig:sphericalintro}. So that each point in space that does not lie on the $z$-axis is defined uniquely, we will restrict $\rho \geq 0$, $0 \leq \theta < 2\pi$ and $0 \leq \varphi \leq \pi$.

The following Key Idea gives conversions to/from our three spatial coordinate systems.

\mnote[1.5in]{\textbf{Note:} The symbol $\rho$ is the Greek letter ``rho.'' Traditionally it is used in the spherical coordinate system, while $r$ is used in the polar and cylindrical coordinate systems.% In space, $r$ and $\rho$ play different roles. In spherical coordinates, $\rho$ is the dist
}

{\tcbset{grow to right by=2em}
\begin{keyidea}[Converting Between Rectangular, Cylindrical and\\Spherical Coordinates]\label{idea:convert_coords}%
\mbox{}\\[-\baselineskip]\parbox[t]{\linewidth}{%
\begin{description}
\item[Rectangular and Cylindrical]
\begin{align*}
r^2 &= x^2+y^2, & \tan \theta &= y/x, & z&=z\\
x&=r\cos \theta, & y&=r\sin\theta, & z&=z
\end{align*}

\item[Rectangular and Spherical]
\begin{align*}
\rho &= \sqrt{x^2+y^2+z^2}, & \tan \theta &= \frac yx, & \cos \varphi &= \frac z{\sqrt{x^2+y^2+z^2}}\\
x&=\rho\sin\varphi\cos\theta, & y&=\rho\sin\varphi\sin\theta, & z&=\rho\cos\varphi
\end{align*}

\item[Cylindrical and Spherical]
\begin{align*}
\rho&=\sqrt{r^2+z^2}, & \theta &= \theta,& \tan \varphi &= r/z \\
r&=\rho \sin \varphi, & \theta &= \theta, & z&=\rho\cos\varphi
\end{align*}
\end{description}}
\end{keyidea}}

\mnote{\textbf{Note:} The role of $\theta$ and $\varphi$ in spherical coordinates differs between mathematicians and physicists. When reading about physics in spherical coordinates, be careful to note how that particular author uses these variables and recognize that these identities may be rearranged.}

\begin{example}[Converting between rectangular and spherical coordinates]\label{ex_spherical4}%
Convert the rectangular point $(2,-2,1)$ to spherical coordinates, and convert the spherical point $(6,\pi/3,\pi/2)$ to rectangular and cylindrical coordinates.
\solution
This rectangular point is the same as used in \autoref{ex_cylindrical4}. Using \autoref{idea:convert_coords}, we find $\rho = \sqrt{2^2+(-1)^2+1^2} = 3$. Using the same logic as in \autoref{ex_cylindrical4}, we find $\theta = 7\pi/4$. Finally, $\cos\varphi = 1/3$, giving $\varphi = \cos^{-1}(1/3) \approx 1.23$, or about $70.53^\circ$. Thus the spherical coordinates are approximately $(3,7\pi/4,1.23)$.

Converting the spherical point $(6,\pi/3,\pi/2)$ to rectangular, we have
\begin{align*}
 x &= 6\sin(\pi/2)\cos(\pi/3) = 3,\\
 y &= 6\sin(\pi/2)\sin(\pi/3) = 3\sqrt{3},\text{ and}\\
 z &= 6\cos(\pi/2) = 0.
\end{align*}
Thus the rectangular coordinates are $(3,3\sqrt{3},0)$.

To convert this spherical point to cylindrical, we have $r = 6\sin(\pi/2) = 6$, $\theta = \pi/3$ and $z = 6\cos(\pi/2) =0$, giving the cylindrical point $(6,\pi/3,0)$.
\end{example}

\begin{example}[Canonical surfaces in spherical coordinates]\label{ex_spherical1}%
Describe the surfaces $\rho=1$, $\theta = \pi/3$ and $\varphi = \pi/6$, given in spherical coordinates.
%
\mtable{Graphing the canonical surfaces in spherical coordinates from \autoref{ex_spherical1}.}{fig:spherical1}{%
\myincludeasythree{
3Droll=0,
3Dortho=0.004559914115816355,
3Dc2c=0.8784526586532593 0.3397093713283539 0.3360334038734436,
3Dcoo=7.126433849334717 10.647956848144531 4.059844017028809,
3Droo=300,
3Dlights=Headlamp}{\marginparwidth}{The 𝜌=1 surface is a sphere centered at the origin. The 𝜃=π/3 surface is the same as with cylindrical coordinates. The 𝜑=π/6 surface is a cone with its point at the origin.}{figures/figspherical1_3D}}
%
\solution
The equation $\rho = 1$ describes all points in space that are 1 unit away from the origin: this is the sphere of radius 1, centered at the origin.

The equation $\theta = \pi/3$ describes the same surface in spherical coordinates as it does in cylindrical coordinates: beginning with the line $\theta = \pi/3$ in the $x$-$y$ plane as given by polar coordinates, extend the line parallel to the $z$-axis, forming a plane.

The equation $\varphi=\pi/6$ describes all points $P$ in space where the ray from the origin to $P$ makes an angle of $\pi/6$ with the positive $z$-axis. This describes a cone, with the positive $z$-axis its axis of symmetry, with point at the origin.

All three surfaces are graphed in \autoref{fig:spherical1}. Note how their intersection uniquely defines the point $P=(1,\pi/3,\pi/6)$.
\end{example}

Spherical coordinates are useful when describing certain domains in space, allowing us to evaluate triple integrals over these domains more easily than if we used rectangular coordinates or cylindrical coordinates. The crux of setting up a triple integral in spherical coordinates is appropriately describing the ``small amount of volume,'' $dV$, used in the integral.

\mtable{Approximating the volume of a standard region in space using spherical coordinates.}{fig:sphericalwedge}{%
\myincludeasythree{
3Droll=0,
3Dortho=0.004123700316995382,
3Dc2c=0.8697938919067383 0.43597298860549927 0.23105460405349731,
3Dcoo=-19.780941009521484 32.02916717529297 43.93011474609375,
3Droo=300.000013020833,
3Dlights=Headlamp}{\marginparwidth}{A region in space near the spherical point (𝜌,𝜃,𝜑)  where each coordinate changes by a small amount, Δ𝜌, Δ𝜃, and Δ𝜑.}{figures/figsphericalwedge_3D}}

Considering \autoref{fig:sphericalwedge}, we can make a small ``spherical wedge'' by varying $\rho$, $\theta$ and $\varphi$ each a small amount, $\Delta\rho$, $\Delta\theta$ and $\Delta\varphi$, respectively. This wedge is approximately a rectangular solid when the change in each coordinate is small, giving a volume of about
\[
\Delta V \approx \Delta\rho\ \times\ \rho\Delta\varphi\ \times\ \rho\sin(\varphi)\Delta\theta.
\]

Given a region $D$ in space, we can approximate the volume of $D$ with many such wedges. As the size of each of $\Delta\rho$, $\Delta\theta$ and $\Delta\varphi$ goes to zero, the number of wedges increases to infinity and the volume of $D$ is more accurately approximated, giving
\[
\dd V = \dd\rho\ \times\ \rho\dd\varphi\ \times\ \rho\sin(\varphi)\dd\theta
=\rho^2\sin(\varphi)\dd\rho\dd\theta\dd\varphi.
\]

Again, this development of $\dd V$ should sound reasonable, and the following theorem states it is the appropriate manner by which triple integrals are to be evaluated in spherical coordinates.

\mnote[3\baselineskip]{\textbf{Note:} It is often most intuitive to evaluate the triple integral in \autoref{thm:triple_int_spherical} by integrating with respect to $\rho$ first; it often does not matter whether we next integrate with respect to $\theta$ or $\varphi$. Different texts present different standard orders, some preferring $\dd\varphi\dd\theta$ instead of $\dd\theta\dd\varphi$. As the bounds for these variables are usually constants in practice, it generally is a matter of preference.}

\begin{theorem}[Triple Integration in Spherical Coordinates]\label{thm:triple_int_spherical}%
Let $w=h(\rho,\theta,\varphi)$ be a continuous function on a closed, bounded region $D$ in space, bounded in spherical coordinates by $\alpha_1 \leq \varphi \leq \alpha_2$, $\beta_1 \leq \theta \leq \beta_2$ and $f_1(\theta,\varphi) \leq \rho \leq f_2(\theta,\varphi)$. Then \index{integration!with spherical coordinates}
\[
\iiint_D h(\rho,\theta,\varphi)\dd V = \int_{\alpha_1}^{\alpha_2}\int_{\beta_1}^{\beta_2}\int_{f_1(\theta,\varphi)}^{f_2(\theta,\varphi)} h(\rho,\theta,\varphi) \rho^2\sin(\varphi)\dd\rho\dd\theta\dd\varphi.
\]
\end{theorem}

\begin{example}[Establishing the volume of a sphere]\label{ex_spherical2}%
Let $D$ be the region in space bounded by the sphere, centered at the origin, of radius $r$. Use a triple integral in spherical coordinates to find the volume $V$ of $D$.
\solution
The sphere of radius $r$, centered at the origin, has equation $\rho = r$. To obtain the full sphere, the bounds on $\theta$ and $\varphi$ are $0\leq \theta \leq 2\pi$ and $0 \leq \varphi \leq \pi$. This leads us to:
\begin{align*}
V &= \iiint_D\dd V\\
	&= \int_0^{\pi}\int_0^{2\pi}\int_0^r\bigl(\rho^2\sin(\varphi)\bigr)\dd\rho\dd\theta\dd\varphi\\
	&= \int_0^\pi\int_0^{2\pi}\left(\frac13\rho^3\sin(\varphi)\Bigr|_{\rho=0}^{\rho=r}\right)\dd\theta\dd\varphi\\
	&= \int_0^\pi\int_0^{2\pi} \left(\frac13r^3\sin(\varphi)\right)\dd\theta\dd\varphi\\
	&= \int_0^\pi \left(\frac{2\pi}3r^3\sin(\varphi)\right)\dd\varphi\\
	&= \left.\left(-\frac{2\pi}3r^3\cos(\varphi)\right)\right|_{\varphi=0}^{\varphi=\pi}\\
	&= \frac{4\pi}3r^3,
\end{align*}
the familiar formula for the volume of a sphere. Note how the integration steps were easy, not using square-roots nor integration steps such as Substitution.
\end{example}

\mtable{Graphing the solid, and its center of mass, from \autoref{ex_spherical3}.}{fig:spherical3}{%
\myincludeasythree{
3Droll=0,
3Dortho=0.004498619586229324,
3Dc2c=0.7278125882148743 0.6620227098464966 0.17892660200595856,
3Dcoo=-23.626569747924805 -17.348602294921875 60.14314270019531,
3Droo=150,
3Dlights=Headlamp}{\marginparwidth}{A cone centered on the positive z axis with its tip at the origin.  The top of the cone is covered by a spherical cap.}{figures/figspherical3_3D}}

\begin{example}[Finding the center of mass using spherical coordinates]\label{ex_spherical3}%
Find the center of mass of the solid with constant density enclosed above by $\rho=4$ and below by $\varphi = \pi/6$, as illustrated in \autoref{fig:spherical3}.
\solution
We will set up the four triple integrals needed to find the center of mass (i.e., to compute $M$, $M_{yz}$, $M_{xz}$ and $M_{xy}$) and leave it to the reader to evaluate each integral. Because of symmetry, we expect the $x$- and $y$- coordinates of the center of mass to be 0.

While the surfaces describing the solid are given in the statement of the problem, to describe the full solid $D$, we use the following bounds: $0 \leq \rho \leq 4$, $0 \leq \theta \leq 2\pi$ and $0 \leq \varphi \leq \pi/6$. Since density $\delta$ is constant, we assume $\delta =1$.

The mass of the solid:
\begin{align*}
M &= \iiint_D\dd m = \iiint_D\dd V\\
	&= \int_0^{\pi/6}\int_0^{2\pi}\int_0^4\bigl(\rho^2\sin(\varphi)\bigr)\dd\rho\dd\theta\dd\varphi\\
	&= \frac{64}3\bigl(2-\sqrt{3}\bigr)\pi \approx 17.958.
\end{align*}

To compute $M_{yz}$, the integrand is $x$; using \autoref{idea:convert_coords}, we have $x = \rho\sin\varphi\cos\theta$. This gives:
\begin{align*}
M_{yz} &= \iiint_D x\dd m \\
	&= \int_0^{\pi/6}\int_0^{2\pi}\int_0^4 \bigl((\rho\sin(\varphi)\cos(\theta))\rho^2\sin(\varphi)\bigr) \dd\rho\dd\theta\dd\varphi\\
	&= \int_0^{\pi/6}\int_0^{2\pi}\int_0^4 \bigl(\rho^3\sin^2(\varphi)\cos(\theta)\bigr) \dd\rho\dd\theta\dd\varphi\\
	&=0,
\end{align*}
which we expected as we expect $\overline{x} = 0$.

To compute $M_{xz}$, the integrand is $y$; using \autoref{idea:convert_coords}, we have $y = \rho\sin\varphi\sin\theta$. This gives:
\begin{align*}
M_{xz} &= \iiint_D y\dd m \\
	&= \int_0^{\pi/6}\int_0^{2\pi}\int_0^4 \bigl((\rho\sin(\varphi)\sin(\theta))\rho^2\sin(\varphi)\bigr) \dd\rho\dd\theta\dd\varphi\\
	&= \int_0^{\pi/6}\int_0^{2\pi}\int_0^4 \bigl(\rho^3\sin^2(\varphi)\sin(\theta)\bigr) \dd\rho\dd\theta\dd\varphi\\
	&=0,
\end{align*}
which we also expected as we expect $\overline{y} = 0$.

To compute $M_{xy}$, the integrand is $z$; using \autoref{idea:convert_coords}, we have $z = \rho\cos\varphi$. This gives:
\begin{align*}
M_{xy} &= \iiint_D z\dd m \\
	&= \int_0^{\pi/6}\int_0^{2\pi}\int_0^4 \bigl((\rho\cos(\varphi))\rho^2\sin(\varphi)\bigr) \dd\rho\dd\theta\dd\varphi\\
	&= \int_0^{\pi/6}\int_0^{2\pi}\int_0^4 \bigl(\rho^3\cos(\varphi)\sin(\varphi)\bigr) \dd\rho\dd\theta\dd\varphi\\
	&=16\pi \approx 50.266.
\end{align*}

Thus the center of mass is $(0,0,M_{xy}/M)=(0,0,3(2+\sqrt3)/4) \approx (0,0,2.799)$, as indicated in \autoref{fig:spherical3}.
\end{example}

% this was originally Schoolcraft / Mecmath

\subsection{Change of Variables in Multiple Integrals}
%\label{sec:jacobian}

Given the difficulty of evaluating multiple integrals, the reader may be wondering if it is possible to simplify those integrals using a suitable substitution for the variables. The answer is yes, though it is a bit more complicated than the substitution method which you learned in single-variable calculus.\index{change of variable}

Recall that if you are given, for example, the definite integral
\[\int_1^2 x^3 \sqrt{x^2 - 1}\dd x ,\]
then you would make the substitution
\begin{align*}
u &= x^2 - 1 \Rightarrow x^2 = u + 1\\
\dd u &= 2x\dd x\\
\intertext{which changes the limits of integration}
x &= 1 \Rightarrow u = 0\\
x &= 2 \Rightarrow u = 3
\end{align*}
so that we get
\begin{align*}
 \int_1^2 x^3 \sqrt{x^2 - 1}\dd x &= \int_1^2 \frac12 x^2 \cdot 2x \sqrt{x^2 - 1}\dd x\\
 &= \int_0^3 \frac12 (u+1)\sqrt{u}\dd u\\
 &= \frac12 \int_0^3 \left( u^{3/2} + u^{1/2} \right)\dd u \\
 &= \frac{14\sqrt3}5 .
\end{align*}
Let us take a different look at what happened when we did that substitution, which will give some motivation for how substitution works in multiple integrals. First, we let $u = x^2 - 1$. On the interval of integration $[ 1,2 ]$, the function $x \mapsto x^2 - 1$ is strictly increasing (and maps $[ 1,2 ]$ onto $[ 0,3 ]$) and hence has an inverse function (defined on the interval $[ 0,3 ]$). That is, on $[ 0,3 ]$ we can define $x$ as a function of $u$, namely
\[x = g(u) = \sqrt{u+1} .\]
Then substituting that expression for $x$ into the function $f(x) = x^3 \sqrt{x^2 - 1}$ gives
\[f(x) = f(g(u)) = (u+1)^{3/2} \sqrt{u} ,\]
and we see that
\begin{align*}
 \frac{\dd x}{\dd u} = g\primeskip'(u) \Rightarrow \dd x &= g\primeskip'(u)\dd u\\
 \dd x &= \frac12 (u+1)^{-1/2}\dd u ,
\end{align*}
so since
\begin{align*}
 g(0) = 1 \Rightarrow 0 = g^{-1} (1)\\
 g(3) = 2 \Rightarrow 3 = g^{-1} (2)\\
\end{align*}
then performing the substitution as we did earlier gives
\begin{align*}
 \int_1^2 f(x)\dd x &= \int_1^2 x^3 \sqrt{x^2 - 1}\dd x\\
 &= \int_0^3 \frac12(u+1)\sqrt{u}\dd u ,\text{ which can be written as}\\
 &= \int_0^3 (u+1)^{3/2} \sqrt{u} \, \cdot \frac12 (u+1)^{-1/2}\dd u ,\text{ which means}\\
 \int_1^2 f(x)\dd x &= \int_{g^{-1}(1)}^{g^{-1}(2)} f(g(u))\,g\primeskip'(u)\dd u .
\end{align*}

In general, if $x = g(u)$ is a one-to-one, differentiable function from an interval $[ c,d ]$ (which you can think of as being on the ``$u$-axis'') onto an interval $[ a,b ]$ (on the $x$-axis), which means that $g\primeskip'(u) \ne 0$ on the interval $(c,d)$, so that $a=g(c)$ and $b=g(d)$, then $c=g^{-1}(a)$ and $d=g^{-1}(b)$, and
\[\int_a^b f(x)\dd x = \int_{g^{-1}(a)}^{g^{-1}(b)} f(g(u))\,g\primeskip'(u)\dd u .\]
This is called the \emph{change of variable} formula for integrals of single-variable functions, and it is what you were implicitly using when doing integration by substitution. This formula turns out to be a special case of a more general formula which can be used to evaluate multiple integrals. We will state the formulas for double and triple integrals involving real-valued functions of two and three variables, respectively. We will assume that all the functions involved are continuously differentiable and that the regions and solids involved all have ``reasonable'' boundaries. The proof of the following theorem is beyond the scope of the text.% See \cite[\S\,15.32 and \S\,15.62]{tm} for all the details.

{\tcbset{grow to right by=2em}
\begin{theorem}[Change of Variables Formula for Multiple Integrals]\label{thm:changevarint}%
Let $x=x(u,v)$ and $y=y(u,v)$ define a one-to-one mapping of a region $R'$ in the $uv$-plane onto a region $R$ in the $xy$-plane such that the determinant\index{change of variable}
  \begin{equation}\label{eqn:jacob2}%
   J(u,v) =
    \begin{vmatrix}
     \dfrac{\partial x}{\partial u} & \dfrac{\partial x}{\partial v}\vspace{2mm}\\
     \dfrac{\partial y}{\partial u} & \dfrac{\partial y}{\partial v}
    \end{vmatrix}
  \end{equation}
  is never $0$ in $R'$. Then
  \begin{equation}\label{eqn:changevarint2}%
   \iint_{R} f(x,y)\dd A(x,y)
   = \iint_{R'} f(x(u,v),y(u,v))\,\abs{J(u,v)}\dd A(u,v) .
  \end{equation}
  We use the notation $\dd A(x,y)$ and $\dd A(u,v)$ to denote the area element in the $(x,y)$ and $(u,v)$ coordinates, respectively.

  Similarly, if $x=x(u,v,w)$, $y=y(u,v,w)$ and $z=z(u,v,w)$ define a one-to-one mapping of a solid $S'$ in $uvw$-space onto a solid $S$ in $xyz$-space such that the determinant
  \begin{equation}\label{eqn:jacob3}%
   J(u,v,w) =
    \begin{vmatrix}
     \dfrac{\partial x}{\partial u} & \dfrac{\partial x}{\partial v} & \dfrac{\partial x}{\partial w}\vspace{2mm}\\
     \dfrac{\partial y}{\partial u} & \dfrac{\partial y}{\partial v} & \dfrac{\partial y}{\partial w}\vspace{2mm}\\
     \dfrac{\partial z}{\partial u} & \dfrac{\partial z}{\partial v} & \dfrac{\partial z}{\partial w}
    \end{vmatrix}
  \end{equation}
  is never $0$ in $S'$, then
  \begin{multline}\label{eqn:changevarint3}%
   \iiint_{S} f(x,y,z)\dd V(x,y,z) =\\
   \iiint_{S'} f(x(u,v,w),y(u,v,w),z(u,v,w))\,\abs{J(u,v,w)}\dd V(u,v,w).
  \end{multline}
\end{theorem}}

\mnote{There are two different uses of $\abs{\cdot}$ in these equations.  The first $\abs{\cdot}$ refers to the Jacobian, which is the determinant of the matrix.  We then take the absolute value of this determinant.}

The determinant $J(u,v)$ in \autoeqref{eqn:jacob2} is called the \textbf{Jacobian} of $x$ and $y$ with respect to $u$ and $v$, and is sometimes written as\index{Jacobian}
\[J(u,v) = \abs{\frac{\partial (x,y)}{\partial (u,v)}} .\]
Similarly, the Jacobian $J(u,v,w)$ of three variables is sometimes written as
\[J(u,v,w) = \abs{\frac{\partial (x,y,z)}{\partial (u,v,w)}} .\]
Notice that \autoeqref{eqn:changevarint2} is saying that $\dd A(x,y) = \abs{J(u,v)}\dd A(u,v)$, which you can think of as a two-variable version of the relation $\dd x = g\primeskip'(u)\dd u$ in the single-variable case.

\youtubeVideo{Bw5yEqwMjQU}{Jacobian}

The following example shows how the change of variables formula is used.

\begin{example}[Change of Variables]\label{ex_simple_Jacobian}%
Evaluate $\ds\iint_{R} e^{\frac{x-y}{x+y}} \dd A$, where $R= \{(x,y): x \ge 0, y \ge 0, x + y \le 1 \}$.
\solution
First, note that evaluating this double integral \emph{without} using substitution is probably impossible, at least in a closed form. By looking at the numerator and denominator of the exponent of $e$, we will try the substitution $u=x-y$ and $v=x+y$. To use the change of variables \autoeqref{eqn:changevarint2}, we need to write both $x$ and $y$ in terms of $u$ and $v$. So solving for $x$ and $y$ gives $x=\frac12(u+v)$ and $y=\frac12(v-u)$. In \autoref{fig:exyint} below, we see how the mapping $x=x(u,v)=\frac12(u+v)$,
 $y=y(u,v)=\frac12(v-u)$ maps the region $R'$ onto $R$ in a one-to-one manner.
 
\noindent\begin{minipage}[t]{\linewidth}\tagstructbegin{tag=Sect}\noindent%
\captionsetup{type=figure}%
 \centering
  \begin{tikzpicture}[alt={On the left is a triangle labeled R in the xy plane with vertices at the origin, (1,0), and (0,1).  The diagonal is x+y=1.  On the right is a triangle labeled R' in the uv plane with vertices at the origin and (±1,1).  An arrow points from right to left with x=(u+v)/2 and y=(v-u)/2.},scale=2]
   % the left area
   \draw[draw={\colorone},fill={\coloronefill},thick]
    (0,0)--(1,0)--node[pos=.5,above right]{\scriptsize$x+y=1$}(0,1)--cycle;
   \draw[draw={\colorone},->](-.2,0)
    --node[above,color=black,pos=1]{\scriptsize$x$}(1.2,0);
   \node[below]at(1,0){\scriptsize$1$};
   \draw[draw={\colorone},->](0,-.2)
    --node[right,color=black,pos=1]{\scriptsize$y$}(0,1.2);
   \node[left]at(0,1){\scriptsize$1$};
   \node[below right] at (0,0) {\scriptsize$0$};
   \node at (0.2,0.2) {\small $R$};
   % the arrow
   \draw[->,very thick](2.5,.7)
    --node[pos=.5,align=center]{$x=\frac12(u+v)$\\[1ex]$y=\frac12(v-u)$}(1.5,.7);
   % the right area
   \draw[draw={\colorone},fill={\coloronefill},thick](4,0)
    --node[pos=.5,below right,color=black]{\scriptsize$u=v$}(5,1)--(3,1)
    --node[pos=.5,below left ,color=black]{\scriptsize$u=-v$}(4,0);
   \draw[draw={\colorone},->](2.8,0)
    --node[above,color=black,pos=1]{\scriptsize$u$}(5.2,0);
   \node[below]at(3,0){\scriptsize$-1$};
   \node[below]at(5,0){\scriptsize$1$};
   \draw[draw={\colorone},->](4,-.2)
    --node[right,color=black,pos=1]{\scriptsize$v$}(4,1.2);
   \node[above left]at(4,1){\scriptsize$1$};
   \node[right]at(4,.6){$R'$};
  \end{tikzpicture}
 \caption{The regions $R$ and $R'$}
 \label{fig:exyint}\tagstructend
\end{minipage}

 Now we see that
 \[
 J(u,v) =
   \begin{vmatrix}
    \dfrac{\partial x}{\partial u} & \dfrac{\partial x}{\partial v}\vspace{2mm}\\
    \dfrac{\partial y}{\partial u} & \dfrac{\partial y}{\partial v}
   \end{vmatrix}
   =\begin{vmatrix}
    \frac12 & \frac12\vspace{2mm}\\
    -\frac12 & \frac12
   \end{vmatrix}
   = \frac12 \Rightarrow \abs{J(u,v)} = \abs{\frac12} = \frac12 ,
 \]
 so using horizontal slices in $R'$, we have
 \begin{align*}
  \iint_{R} e^{\frac{x-y}{x+y}} \dd A &= \iint_{R'} f(x(u,v),y(u,v))\,\abs{J(u,v)}\dd A\\
  &= \int_0^1 \int_{-v}^v e^{\frac uv} \, \frac12\dd u\dd v\\
  &= \int_0^1 \left(\left. \frac v2 e^{\frac uv} \,\right|_{u=-v}^{u=v} \right) \dd v\\
  &= \int_0^1 \frac v2 (e - e^{-1})\dd v\\
  &= \left.\frac{v^2}4 (e - e^{-1}) \,\right|_{0}^{1}
  = \frac14 \left( e - \frac1e \right) = \frac{e^2 - 1}{4e}.
 \end{align*}
\end{example}

The change of variables formula can be used to evaluate double integrals in polar coordinates. Letting
\[
 x = x(r,\theta) = r\cos \theta \quad \text{and} \quad y = y(r,\theta) = r\sin \theta ,
\]
we have
\begin{multline*}
 J(u,v) =
  \begin{vmatrix}
   \dfrac{\partial x}{\partial r} & \dfrac{\partial x}{\partial \theta}\vspace{2mm}\\
   \dfrac{\partial y}{\partial r} & \dfrac{\partial y}{\partial \theta}
  \end{vmatrix} =
  \begin{vmatrix}
   \cos \theta & -r\sin \theta\vspace{2mm}\\
   \sin \theta & r\cos \theta
  \end{vmatrix}= r\cos^2 \theta + r\sin^2 \theta = r \\
  \Rightarrow \abs{J(u,v)} = \abs{r} = r ,
\end{multline*}
which verifies \autoref{idea:doublepol}.
%so we have the following formula:
%
%\begin{keyidea}[Double Integral in Polar Coordinates]\label{idea_polar_int}%
% \begin{equation}\label{eqn:polarint}%
%  \iint_{R} f(x,y)\dd x\dd y
%  = \iint_{R'} f(r\cos \theta,r\sin \theta)\,r\,dr\,d\theta ,
% \end{equation}\index{double integral!polar coordinates}\index{coordinates!polar}
% where the mapping $x=r\cos \theta$, $y=r\sin \theta$ maps the region $R'$ in the $r\theta$-plane onto the region $R$ in the $xy$-plane in a one-to-one manner.
%\end{keyidea}
%
%\begin{example}[Integrating in Polar Coordinates]\label{ex_polar_parab}%
%Find the volume $V$ inside the paraboloid $z=x^2 + y^2$ for $0 \le z \le 1$.
%\solution
%Using vertical slices, we see that
%\[V = \iint_{R} (1 - z)\dd A = \iint_{R} (1 - (x^2 + y^2 ))\dd A ,\]
%where $R = \lbrace (x,y): x^2 + y^2 \le 1 \rbrace$ is the unit disk in $\mathbb{R}^2$ (see \autoref{fig_parab}). In polar coordinates $(r,\theta)$ we know that $x^2 + y^2 = r^2$ and that the unit disk $R$ is the set $R' = \lbrace (r,\theta):0 \le r \le 1, 0 \le \theta \le 2\pi \rbrace$. Thus,
%\mtable{$z=x^2 + y^2$ for \autoref{ex_polar_parab}}{fig_parab}{\begin{tikzpicture}[alt={A paraboloid cup with vertex at the origin.  It stops at z=x²+y²=1.}]
%  \usetikzlibrary{arrows}
%  \definecolor{insideo}{HTML}{798084}
%  \definecolor{insidei}{HTML}{F0F0F0}
%  \definecolor{outer}{HTML}{424296}
%  \definecolor{inner}{HTML}{D8D8FF}
%  \shadedraw [left color=insideo,right color=insideo,middle color=insidei] (0,3) ellipse (2 and 0.7);
%  \shadedraw [left color=outer,right color=outer,middle color=inner]
%  (2,3) arc (360:180:2 and 0.7) -- (-2,3) parabola bend (0,0) (2,3);
%  \draw [dashed,line width=0.2pt] (-2,3) -- (2,3);
%  \draw [dashed,line width=0.2pt] (-0.66,2.34) -- (0.66,3.66);
%  \draw [line width=0.2pt] (-0.66,2.34) parabola [bend at end] (0,0);
%  \draw [dashed,line width=0.2pt] (0.66,3.66) parabola [bend at end] (0,0);
%  \draw [line width=0.2pt] (-1.25,1.2) arc (180:360:1.25 and 0.4);
%  \draw [dashed,line width=0.2pt] (-1.25,1.2) arc (180:0:1.25 and 0.4);
%  \draw [black!60,line width=0.3pt,-latex] (-2,0) -- (2,0,0);
%  \draw [black!60,line width=0.3pt,-latex] (0,0) -- (0,4,0);
%  \draw [black!60,line width=0.3pt,-latex] (0,0) -- (0,0,2);
%  \node[above]at(2,0){\small$y$};
%  \node[right]at(0,4){\small$z$};
%  \node[right]at(0,0,2){\small$x$};
%  \node[below right]at(0,0){\small$0$};
%%  \pgfputat{\pgfpointxyz{1.9}{0.2}{0}}{\pgfbox[center,center]{\small $y$}};
%%  \pgfputat{\pgfpointxyz{0.2}{3.9}{0}}{\pgfbox[center,center]{\small $z$}};
%%  \pgfputat{\pgfpointxyz{0.2}{0}{1.8}}{\pgfbox[center,center]{\small $x$}};
%%  \pgfputat{\pgfpointxyz{0.05}{-0.2}{0}}{\pgfbox[center,center]{\small $0$}};
%  \node [right] at (1.5,3.7) {\small $x^2 + y^2 = 1$};
%  \node [above left] at (0,3) {\small $1$};
% \end{tikzpicture}}%
% \begin{align*}
%  V &= \int_0^{2\pi} \int_0^1 (1 - r^2 )\,r\,dr\,d\theta\\
%   &= \int_0^{2\pi} \int_0^1 (r - r^3 )\,dr\,d\theta\\
%   &= \int_0^{2\pi} \left( \left.\frac{r^2}2 - \frac{r^4}4 \,\right|_{r=0}^{r=1} \right) \,d\theta\\
%   &= \int_0^{2\pi} \frac14 \,d\theta\\
%   &= \frac\pi2.
% \end{align*}
%\end{example}
%
%\begin{example}[Integrating in Polar Coordinates]\label{ex_polar_cone}%
%Find the volume $V$ inside the cone $z=\sqrt{x^2 + y^2}$ for $0 \le z \le 1$.
%\solution
% \mtable{$z=\sqrt{x^2 + y^2}$ for \autoref{ex_polar_cone}}{fig_cone}{\begin{tikzpicture}[alt={A cone with vertex at the origin centered along the positive z axis.  The cone stops at z=x²+y²+1.}]
%  \usetikzlibrary{arrows}
%  \definecolor{insideo}{HTML}{798084}
%  \definecolor{insidei}{HTML}{F0F0F0}
%  \definecolor{outer}{HTML}{424296}
%  \definecolor{inner}{HTML}{D8D8FF}
%  \shadedraw [left color=insideo,right color=insideo,middle color=insidei] (1.3,2) arc (0:180:1.3 and .5) --
%   (-1.3,2) arc (180:360:1.3 and .5);
%  \shadedraw [left color=outer,right color=outer,middle color=inner] (0,0) -- (-1.3,2) arc (180:360:1.3 and .5)
%   -- (0,0);
%  \draw [dashed,line width=0.2pt] (-1.3,2) -- (1.3,2);
%  \draw [black!60,line width=0.3pt,-latex] (-2,0) -- (2,0,0);
%  \draw [black!60,line width=0.3pt,-latex] (1,1) -- (0,0,2);
%  \draw [black!60,line width=0.3pt,-latex] (0,-0.5) -- (0,2.9,0);
%  \node[above]at(2,0){\small$y$};
%  \node[right]at(0,3){\small$z$};
%  \node[right]at(0,0,2){\small$x$};
%  \node[below right]at(0,0){\small$0$};
%%  \pgfputat{\pgfpointxyz{1.9}{0.2}{0}}{\pgfbox[center,center]{\small $y$}};
%%  \pgfputat{\pgfpointxyz{0.2}{2.8}{0}}{\pgfbox[center,center]{\small $z$}};
%%  \pgfputat{\pgfpointxyz{0.2}{0}{1.8}}{\pgfbox[center,center]{\small $x$}};
%%  \pgfputat{\pgfpointxyz{0.15}{-0.2}{0}}{\pgfbox[center,center]{\small $0$}};
%  \node [above] at (1.3,2.4) {\small $x^2 + y^2 = 1$};
%  \node [above left] at (0,2) {\small $1$};
% \end{tikzpicture}
% }
% Using vertical slices, we see that
% \[
%  V = \iint_{R} (1 - z)\dd A
%  = \iint_{R} \left( 1 - \sqrt{x^2 + y^2} \right) \dd A ,
% \]
% where $R = \lbrace (x,y): x^2 + y^2 \le 1 \rbrace$ is the unit disk in $\mathbb{R}^2$ (see \autoref{fig_cone}). In polar coordinates $(r,\theta)$ we know\\that $\sqrt{x^2 + y^2} = r$ and that the unit disk $R$ is the set\\$R' = \lbrace (r,\theta):0 \le r \le 1, 0 \le \theta \le 2\pi \rbrace$. Thus,
% \begin{align*}
%  V &= \int_0^{2\pi} \int_0^1 (1 - r)\,r\,dr\,d\theta\\
%   &= \int_0^{2\pi} \int_0^1 (r - r^2 )\,dr\,d\theta\\
%   &= \int_0^{2\pi} \left( \left.\frac{r^2}2 - \frac{r^3}3 \,\right|_{r=0}^{r=1} \right) \,d\theta\\
%   &= \int_0^{2\pi} \frac16 \,d\theta\\
%   &= \frac\pi3.
% \end{align*}
%\end{example}
%
In a similar fashion, %it can be shown (see
Exercises \ref{pr_cyl_jac} and \ref{pr_sph_jac}
ask you to verify Theorems \ref{thm:triple_int_cylindrical} and \ref{thm:triple_int_spherical} by computing the appropriate Jacobian.\bigskip
%) that triple integrals in cylindrical and spherical coordinates take the following forms:

%\begin{keyidea}[Triple Integral in Cylindrical Coordinates]\label{idea_cylin_int}%
% \begin{equation}\label{eqn:cylinderint}%
%  \iiint_{S} f(x,y,z)\dd x\dd y\dd z = \iiint_{S'} f(r\cos \theta,r\sin \theta,z)\,r\,dr\,d\theta\,dz ,
% \end{equation}\index{triple integral!cylindrical coordinates}
% where the mapping $x=r\cos \theta$, $y=r\sin \theta$, $z=z$ maps the solid $S'$ in  $r\theta z$-space onto the solid $S$ in $xyz$-space in a one-to-one manner.
%\end{keyidea}
%
%\begin{keyidea}[Triple Integral in Spherical Coordinates]\label{idea_sphere_int}%
%\mbox{}\\[-2\baselineskip]
% \begin{multline}\label{eqn:sphereint}%
%  \iiint_S f(x,y,z)\dd x\dd y\dd z =\\
%  \iiint_{S'} f(\rho\sin \phi\,\cos \theta,\rho\sin \phi\,\sin \theta,\rho\cos \phi)
%  \,\rho^2 \,\sin \phi \,d\rho\,d\phi\,d\theta ,
% \end{multline}\index{triple integral!spherical coordinates}
% where the mapping $x=\rho\sin \phi\,\cos \theta$, $y=\rho\sin \phi\,\sin \theta$, $z=\rho\cos \phi$ maps the solid $S'$ in $\rho\phi\theta$-space onto the solid $S$ in $xyz$-space in a one-to-one manner.
%\end{keyidea}
%
%\begin{example}[Triple Integral over a Region in Cylindrical Coordinates]\label{exmp_cyl_cone_octant}%
%Evaluate $\iiint_D x\dd V$ where $D$ is the region bounded by the first octant, the cone $z=\sqrt{x^2+y^2}$, and $z=3$.
%\solution
%In cylindrical coordinates $D$ is given by $0\le\theta\leq\frac\pi2$, $0\leq r\leq3$, and $r\leq z\leq3$.  Now we have
%\begin{align*}
%\iiint_D x\dd V
%&=\int_0^{\frac\pi2}\int_0^3\int_r^3 r\cos\theta r\,dz\,dr\,d\theta \\
%&=\left.\int_0^{\frac\pi2}\int_0^3 r^2\,z\,\cos \theta\, \right|_{z=r}^{z=3}\,dz\,dr\,d\theta \\
%&=\int_0^{\frac\pi2}\int_0^3 (3r^2-r^3) \cos \theta \,dr\,d\theta \\
%&=\left.\int_0^{\frac\pi2} (r^3-\frac14 r^4) \cos \theta \,\right|_{r=0}^{r=3} \,d\theta \\
%&=\int_0^{\frac\pi2} \frac{27}4 \cos \theta \,d\theta \\
%&=\left.\frac{27}4 \sin\theta \,\right|_{\theta=0}^{\theta=\frac\pi2} \\  
%&=\frac{27}4.
%\end{align*}
%\end{example}
%
%\begin{example}[Finding the Volume of a Sphere]\label{exmp_volsphere}%
%For $a > 0$, find the volume $V$ inside the sphere $S = x^2 + y^2 + z^2 = a^2$.
%We see that $S$ is the set $\rho = a$ in spherical coordinates, so
% \begin{align*}
%  V &= \iiint_S 1\dd V
%   = \int_0^{2\pi} \int_0^\pi \int_0^a 1 \,\rho^2 \,\sin \phi \,d\rho\,d\phi\,d\theta\\
%   &= \int_0^{2\pi} \int_0^\pi \left( \left.\frac{\rho^3}3\,\right|_{\rho=0}^{\rho=a}\right)\,\sin \phi \,d\phi\,d\theta
%   = \int_0^{2\pi} \int_0^{\pi} \frac{a^3}3\,\sin \phi\,d\phi\,d\theta\\
%   &= \int_0^{2\pi} \left( \left. -\frac{a^3}3\,\cos \phi\,\right|_{\phi=0}^{\phi=\pi} \right) \,d\theta
%   = \int_0^{2\pi} \frac{2a^3}3\,d\theta
%   = \frac{4\pi a^3}3 .
% \end{align*}
%\end{example}

%This chapter investigated the natural follow--on to partial derivatives: iterated integration. We learned how to use the bounds of a double integral to describe a region in the plane using both rectangular and polar coordinates, then later expanded to use the bounds of a triple integral to describe a region in space. We used double integrals to find volumes under surfaces, surface area, and the center of mass of lamina; we used triple integrals as an alternate method of finding volumes of space regions and also to find the center of mass of a region in space.
%
%Integration does not stop here. We could continue to iterate our integrals, next investigating ``quadruple integrals'' whose bounds describe a region in 4--dimensional space (which are very hard to visualize). We can also look back to ``regular'' integration where we found the area under a curve in the plane. A natural analogue to this is finding the ``area under a curve,'' where the curve is in space, not in a plane. These are just two of many avenues to explore under the heading of ``integration.''

%\printexercises{exercises/13-Jacobian-exercises}


This section has provided a brief introduction into two new coordinate systems useful for identifying points in space. Each can be used to define a variety of surfaces in space beyond the canonical surfaces graphed as each system was introduced.

However, the usefulness of these coordinate systems does not lie in the variety of surfaces that they can describe nor the regions in space these surfaces may enclose. Rather,  cylindrical coordinates are mostly used to describe cylinders and spherical coordinates are mostly used to describe spheres. These shapes are of special interest in the sciences, especially in physics, and computations on/inside these shapes is difficult using rectangular coordinates. For instance, in the study of electricity and magnetism, one often studies the effects of an electrical current passing through a wire; that wire is essentially a cylinder, described well by cylindrical coordinates.\bigskip

This chapter investigated the natural follow-on to partial derivatives: iterated integration. We learned how to use the bounds of a double integral to describe a region in the plane using both rectangular and polar coordinates, then later expanded to use the bounds of a triple integral to describe a region in space. We used double integrals to find volumes under surfaces, surface area, and the center of mass of lamina; we used triple integrals as an alternate method of finding volumes of space regions and also to find the center of mass of a region in space.

Integration does not stop here. We could continue to iterate our integrals, next investigating ``quadruple integrals'' whose bounds describe a region in 4-dimensional space (which are very hard to visualize). We can also look back to ``regular'' integration where we found the area under a curve in the plane. A natural analogue to this is finding the ``area under a curve,'' where the curve is in space, not in a plane. These are just two of many avenues to explore under the heading of ``integration.''

\printexercises{exercises/13-07-exercises}
